% Compilar desde este directorio con:
% pdflatex Clase_05_Diseno_experimental_de_Machine_Learning.tex
\documentclass[aspectratio=169,10pt]{beamer}

\usepackage[utf8]{inputenc}
\usepackage[T1]{fontenc}
\usepackage[spanish,es-nodecimaldot]{babel}
\usepackage{amsmath,amssymb,booktabs,tabularx,array,tikz}
\usetikzlibrary{arrows.meta,positioning,calc,shapes.geometric}

\definecolor{UAPNavy}{HTML}{17324D}
\definecolor{UAPTeal}{HTML}{007F82}
\definecolor{UAPGold}{HTML}{E2A33A}
\definecolor{UAPLight}{HTML}{F3F6F8}
\definecolor{UAPGray}{HTML}{536471}
\definecolor{UAPRed}{HTML}{A94343}

\tikzset{
  flow/.style={-{Stealth[length=2.2mm]},thick,draw=UAPNavy},
  datum/.style={draw=UAPGold,fill=UAPGold!16,rounded corners=2pt,align=center,minimum height=8mm,text=UAPNavy},
  process/.style={draw=UAPTeal,fill=UAPTeal!12,rounded corners=3pt,very thick,align=center,minimum height=10mm,text=UAPNavy},
  decision/.style={draw=UAPNavy,fill=UAPNavy!9,rounded corners=3pt,very thick,align=center,minimum height=10mm,text=UAPNavy},
  warning/.style={draw=UAPRed,fill=UAPRed!7,rounded corners=3pt,align=left},
  axis/.style={-{Stealth[length=2mm]},draw=UAPGray,thick}
}

\usetheme{Boadilla}
\usecolortheme{default}
\setbeamercolor{normal text}{fg=UAPNavy,bg=white}
\setbeamercolor{structure}{fg=UAPTeal}
\setbeamercolor{frametitle}{fg=white,bg=UAPNavy}
\setbeamercolor{title}{fg=white}
\setbeamercolor{subtitle}{fg=UAPGold}
\setbeamercolor{block title}{fg=white,bg=UAPTeal}
\setbeamercolor{block body}{fg=UAPNavy,bg=UAPLight}
\setbeamercolor{alertblock title}{fg=white,bg=UAPRed}
\setbeamercolor{alertblock body}{fg=UAPNavy,bg=UAPRed!7}
\setbeamerfont{frametitle}{series=\bfseries}
\setbeamertemplate{navigation symbols}{}
\setbeamertemplate{itemize item}{\color{UAPGold}$\blacktriangleright$}
\setbeamertemplate{itemize subitem}{\color{UAPTeal}$\bullet$}
\setbeamertemplate{footline}{%
  \leavevmode\hbox{%
    \begin{beamercolorbox}[wd=.78\paperwidth,ht=2.8ex,dp=1.2ex,leftskip=1em]{author in head/foot}%
      \color{white}\insertshorttitle
    \end{beamercolorbox}%
    \begin{beamercolorbox}[wd=.22\paperwidth,ht=2.8ex,dp=1.2ex,rightskip=1em plus1fil]{date in head/foot}%
      \color{UAPGray}\hfill Semana 5 \enspace | \enspace \insertframenumber/\inserttotalframenumber
    \end{beamercolorbox}%
  }%
}

\newcommand{\concepto}[1]{\textcolor{UAPTeal}{\textbf{#1}}}
\newcommand{\br}{\\}
\newcommand{\pregunta}[1]{\begin{block}{Pregunta clave}#1\end{block}}

\title[IA · Clase 5]{Diseño experimental de Machine Learning}
\subtitle{De datos históricos a decisiones confiables}
\author{Curso de Inteligencia Artificial}
\institute{Semana 5 · Clase 5}
\date{}

\begin{document}

% 1
\begin{frame}[plain]
  \begin{beamercolorbox}[wd=\paperwidth,ht=\paperheight,sep=2em]{frametitle}
    \vfill
    {\usebeamerfont{title}\usebeamercolor[fg]{title}\Huge\inserttitle\par}
    \vspace{.45em}{\usebeamerfont{subtitle}\usebeamercolor[fg]{subtitle}\Large\insertsubtitle\par}
    \vspace{1em}
    \begin{tikzpicture}[node distance=4mm,scale=.72,transform shape]
      \node[decision] (d) {Decisión}; \node[datum,right=of d] (x) {Datos};
      \node[process,right=of x] (f) {Función\br aprendida}; \node[datum,right=of f] (e) {Evaluación};
      \node[decision,right=of e] (a) {Acción};
      \draw[flow] (d)--(x); \draw[flow] (x)--(f); \draw[flow] (f)--(e); \draw[flow] (e)--(a);
    \end{tikzpicture}\par
    \vspace{.7em}{\color{white}\small El valor de una predicción depende del experimento que permite confiar en ella.\par}
    \vfill {\color{UAPGold}\large\insertauthor\par}{\color{white}\insertinstitute\par}
  \end{beamercolorbox}
\end{frame}

% 2
\begin{frame}{Continuidad: de incertidumbre a aprendizaje}
  \small Un agente decide con información parcial y costos distintos para cada error.
  \begin{itemize}
    \item La probabilidad representó la creencia sobre un estado oculto.
    \item El costo esperado conectó esa creencia con una acción.
    \item En sistemas reales, las probabilidades y reglas se estiman con datos históricos.
  \end{itemize}
  \begin{block}{Continuidad}Diseñar el experimento determina si una estimación aprendida puede sostener una decisión fuera de los datos usados para construirla.\end{block}
\end{frame}

% 3
\begin{frame}{Propósito y resultados de aprendizaje}
  \small \textbf{Propósito:} diseñar una evaluación de Machine Learning alineada con la decisión y libre de información indebida.
  \begin{block}{Al finalizar, podremos}
  \begin{itemize}
    \item distinguir tarea, modelo, regla de decisión y agente;
    \item definir unidad, población, instante, horizonte, variables y etiqueta;
    \item elegir particiones coherentes con la forma de uso;
    \item construir baselines y métricas vinculadas con costos;
    \item aplicar validación cruzada sin fuga de información;
    \item reconocer subajuste, sobreajuste y comparaciones injustas;
    \item documentar un procedimiento reproducible y congelado.
  \end{itemize}
  \end{block}
\end{frame}

% 4
\begin{frame}{Pregunta de apertura}
  \small
  \pregunta{Un vehículo debe decidir si ingresa a un tramo. ¿Cómo demostramos que su estimación de bloqueo funcionará con tramos futuros?}
  \begin{itemize}
    \item ¿Qué señales existen antes de ingresar? ¿Cuándo se confirma el estado real?
    \item ¿Qué error es más costoso: ingresar con bloqueo o desviarse sin necesidad?
    \item ¿Qué datos representan el uso real?
  \end{itemize}
  \begin{alertblock}{Idea}El experimento debe reproducir la información y las restricciones disponibles al decidir.\end{alertblock}
\end{frame}

% 5
\begin{frame}{De probabilidades dadas a estimaciones aprendidas}
  \small
  \begin{tabularx}{\textwidth}{XX}\toprule
  \textbf{Probabilidad especificada} & \textbf{Estimación aprendida}\\\midrule
  Se entrega $P(B)$ y $P(A\mid B)$ & Se observan ejemplos históricos $(x_i,y_i)$\\
  La tabla define la relación & Una función $f$ aproxima la relación\\
  Bayes actualiza la creencia & $\hat y=f(x)$ estima el objetivo\\\bottomrule
  \end{tabularx}
  \begin{block}{Puente}Los datos reemplazan parámetros dados por estimaciones; el diseño experimental controla si esas estimaciones generalizan.\end{block}
\end{frame}

% 6
\begin{frame}{¿Qué significa aprender?}
  \small Aprender es ajustar, a partir de ejemplos, una función que transforma información disponible en una salida útil:
  \[f:\mathcal{X}\rightarrow\mathcal{Y},\qquad \hat y=f(x).\]
  \begin{itemize}
    \item $x$: variables observadas antes de decidir. \item $y$: resultado real definido por el problema.
    \item $\hat y$: estimación producida por la función aprendida.
    \item El ajuste busca buen desempeño en casos no usados para aprender.
  \end{itemize}
  \begin{block}{Resultado}Memorizar ejemplos no equivale a aprender una relación generalizable.\end{block}
\end{frame}

% 7
\begin{frame}{Modelo, regla de decisión y agente}
  \small
  \centering
  \begin{tikzpicture}[node distance=4mm,scale=.74,transform shape]
    \node[datum] (x) {Señales $x$}; \node[process,right=of x] (f) {Modelo $f$};
    \node[datum,right=of f] (y) {Predicción $\hat y$}; \node[decision,right=of y] (r) {Regla\br y costos};
    \node[datum,right=of r] (a) {Acción}; \draw[flow] (x)--(f); \draw[flow] (f)--(y); \draw[flow] (y)--(r); \draw[flow] (r)--(a);
  \end{tikzpicture}
  \begin{tabularx}{\textwidth}{>{\bfseries}p{.24\textwidth}X}\toprule
  Modelo & estima una clase, valor o puntaje\\ Regla de decisión & convierte salida y costos en una acción\\ Agente & observa, aplica la regla, actúa y recibe consecuencias\\\bottomrule\end{tabularx}
  \begin{alertblock}{Control}Una predicción de bloqueo no ordena desviarse; la acción depende del umbral, los costos y las restricciones.\end{alertblock}
\end{frame}

% 8
\begin{frame}{Anatomía de un ejemplo}
  \small Para un tramo observado a las 08:00:
  \begin{tabularx}{\textwidth}{>{\bfseries}p{.12\textwidth}X}\toprule
  Elemento & Ejemplo\\\midrule
  $x$ & lluvia reciente, velocidad aguas arriba, alertas previas, tipo de vía\\
  $y$ & $1$ si queda confirmado como bloqueado en 15 min; $0$ si queda libre\\
  $\hat y$ & puntaje de bloqueo estimado antes de ingresar\\\bottomrule
  \end{tabularx}
  \begin{block}{Unidad temporal}Una fila representa una unidad en un instante. La etiqueta se conoce después; la predicción usa únicamente el pasado disponible a las 08:00.\end{block}
\end{frame}

% 9
\begin{frame}{Aprendizaje supervisado}
  \small El conjunto de entrenamiento contiene pares $(x_i,y_i)$ con respuesta conocida:
  \[\{(x_1,y_1),\ldots,(x_n,y_n)\}\longrightarrow\hat f.\]
  \begin{itemize}
    \item La etiqueta guía el ajuste de la función.
    \item La salida se compara con respuestas reales no usadas en el ajuste.
    \item Ejemplos: bloqueo de un tramo, demanda de agua y tiempo de resolución de un reclamo.
  \end{itemize}
  \begin{block}{Idea}Supervisado describe la disponibilidad de etiquetas, no un algoritmo particular.\end{block}
\end{frame}

% 10
\begin{frame}{Clasificación y regresión}
  \small
  \begin{tabularx}{\textwidth}{>{\bfseries}p{.22\textwidth}p{.25\textwidth}X}\toprule
  Tarea & Salida & Ejemplo de transferencia\\\midrule
  Clasificación & categoría o puntaje & tramo bloqueado/libre; reclamo urgente/no urgente\\
  Regresión & valor numérico & demanda de agua; minutos para resolver un reclamo\\\bottomrule
  \end{tabularx}
  \begin{itemize}\item En clasificación, un puntaje puede convertirse en clase mediante un umbral.\item En regresión, la salida conserva escala y unidad.\end{itemize}
  \begin{alertblock}{Control}La forma de $y$ define la tarea; la acción pertenece al sistema de decisión.\end{alertblock}
\end{frame}

% 11
\begin{frame}{Aprendizaje no supervisado}
  \small Los datos contienen variables $x_i$, pero no una etiqueta objetivo $y_i$.
  \begin{itemize}
    \item Busca estructura, similitudes o representaciones en los datos.
    \item El \concepto{clustering} agrupa observaciones según un criterio de similitud.
    \item Ejemplo: agrupar patrones de consumo de agua sin categorías previas.
    \item Los grupos obtenidos requieren interpretación y validación respecto del propósito.
  \end{itemize}
  \begin{block}{Resultado}Descubrir estructura no demuestra por sí mismo capacidad predictiva ni prescribe una acción.\end{block}
\end{frame}

% 12
\begin{frame}{Aprendizaje semisupervisado}
  \small Combina pocos ejemplos etiquetados con muchos sin etiqueta:
  \[\mathcal{D}_L=\{(x_i,y_i)\},\quad\mathcal{D}_U=\{x_j\},\quad |\mathcal{D}_L|\ll|\mathcal{D}_U|.\]
  \begin{itemize}
    \item Es útil cuando medir señales es barato y confirmar el resultado es costoso.
    \item En movilidad abundan trazas, pero confirmar cada bloqueo exige verificación.
    \item La evaluación final necesita etiquetas confiables y una partición representativa.
  \end{itemize}
  \begin{alertblock}{Control}Datos sin etiqueta amplían información sobre $x$; no reemplazan una definición precisa de $y$.\end{alertblock}
\end{frame}

% 13
\begin{frame}{Comparación de paradigmas}
  \scriptsize
  \begin{tabularx}{\textwidth}{>{\bfseries}p{.21\textwidth}p{.29\textwidth}X}\toprule
  Paradigma & Evidencia disponible & Pregunta central\\\midrule
  Supervisado & ejemplos $(x,y)$ & ¿qué salida corresponde a un caso nuevo?\\
  No supervisado & ejemplos $x$ & ¿qué estructura aparece en los datos?\\
  Semisupervisado & pocos $(x,y)$ y muchos $x$ & ¿cómo aprovechar ambos conjuntos?\\
  Refuerzo & interacción, acciones y recompensas & ¿qué acciones acumulan mejores consecuencias?\\\bottomrule
  \end{tabularx}
  \begin{block}{Idea}El paradigma surge del tipo de retroalimentación disponible y del objetivo de la tarea.\end{block}
\end{frame}

% 14
\begin{frame}{Partir de la decisión}
  \small El diseño comienza en el uso, no en el archivo de datos:
  \begin{enumerate}
    \item decisión que debe tomar el sistema; \item consecuencias de cada error;
    \item salida predictiva que informa la decisión; \item instante exacto en que se requiere;
    \item evidencia disponible en ese instante; \item evaluación que reproduce esas condiciones.
  \end{enumerate}
  \pregunta{¿Qué cambiará operativamente cuando el modelo produzca $\hat y$?}
\end{frame}

% 15
\begin{frame}{Definición del caso conductor}
  \scriptsize
  \begin{tabularx}{\textwidth}{>{\bfseries}p{.22\textwidth}X}\toprule
  Elemento & Definición\\\midrule
  Decisión & ingresar al tramo o tomar un desvío\\ Salida predictiva & puntaje de bloqueo en los próximos 15 min\\
  Regla & desviar si el puntaje supera el umbral operativo\\ Evidencia & señales registradas antes de ingresar\\
  Confirmación & reporte verificado del estado del tramo\\ Costo principal & seguridad, demora y consumo del recorrido\\\bottomrule
  \end{tabularx}
  \begin{block}{Criterio de consistencia}Cada ejemplo debe respetar estas definiciones para representar la misma tarea.\end{block}
\end{frame}

% 16
\begin{frame}{Unidad y población}
  \small
  \begin{itemize}
    \item \concepto{Unidad de análisis}: un par (tramo, instante de decisión).
    \item \concepto{Población objetivo}: decisiones de ingreso de la flota en la red y condiciones operativas definidas.
    \item \concepto{Muestra}: unidades históricas observadas y recuperadas por el sistema de datos.
  \end{itemize}
  La misma vía en dos instantes produce dos unidades distintas. Dos registros de la misma unidad no son evidencia independiente.
  \begin{alertblock}{Control}Evaluar sobre otra población responde otra pregunta, aunque las columnas coincidan.\end{alertblock}
\end{frame}

% 17
\begin{frame}{Objetivo, etiqueta, predicción y acción}
  \small
  \begin{tabularx}{\textwidth}{>{\bfseries}p{.2\textwidth}X}\toprule
  Objetivo & anticipar bloqueo antes de ingresar\\ Etiqueta $y$ & estado confirmado dentro del horizonte: $1/0$\\
  Predicción $\hat y$ & puntaje o clase estimada\\ Acción & ingresar, desviar o solicitar verificación\\\bottomrule
  \end{tabularx}
  \[x\xrightarrow{f}\hat y\xrightarrow{\text{regla + costos}}a.\]
  \begin{block}{Resultado}Evaluar $\hat y$ mide predicción; evaluar consecuencias de $a$ mide el sistema de decisión.\end{block}
\end{frame}

% 18
\begin{frame}{Instante y horizonte}
  \small
  \begin{itemize}
    \item $t_0$: instante en que el vehículo debe decidir.
    \item Ventana de observación: datos con marca temporal $t\le t_0$.
    \item Horizonte $H=15$ min; etiqueta: bloqueo confirmado en $(t_0,t_0+H]$.
  \end{itemize}
  \centering
  \begin{tikzpicture}[scale=.88]
    \draw[axis] (0,0)--(11,0); \draw[very thick,UAPTeal] (1,0)--(4.5,0); \draw[very thick,UAPGold] (4.5,0)--(8,0);
    \draw (4.5,-.15)--(4.5,.15); \node[above] at (2.7,.1) {pasado observable}; \node[above] at (6.25,.1) {horizonte $H$};
    \node[below,align=center] at (4.5,-.15) {$t_0$: predicción\\ acción}; \node[below] at (8,-.15) {confirmación};
  \end{tikzpicture}
  \begin{alertblock}{Control}Cambiar $t_0$ o $H$ cambia la tarea, la etiqueta y la utilidad operativa.\end{alertblock}
\end{frame}

% 19
\begin{frame}{Variables disponibles al decidir}
  \small
  \begin{tabularx}{\textwidth}{XX}\toprule
  \textbf{Admitidas en $t_0$} & \textbf{No admitidas en $t_0$}\\\midrule
  lluvia acumulada hasta $t_0$ & lluvia registrada después de $t_0$\\
  velocidad aguas arriba hasta $t_0$ & velocidad medida tras ingresar\\
  reportes recibidos antes de $t_0$ & confirmación final del bloqueo\\
  atributos estables del tramo & duración total del incidente\\\bottomrule
  \end{tabularx}
  \begin{alertblock}{Control}Una columna es válida por su momento real de disponibilidad, no por aparecer en la tabla histórica.\end{alertblock}
\end{frame}

% 20
\begin{frame}{Fuga de información}
  \small Existe \concepto{fuga} cuando el entrenamiento o la evaluación reciben información que no estaría disponible al producir la predicción real.
  \begin{itemize}
    \item Infla el resultado medido sin mejorar el uso futuro.
    \item Puede ocurrir en variables, transformaciones o particiones.
    \item También aparece cuando una misma entidad aporta registros casi idénticos a entrenamiento y evaluación.
  \end{itemize}
  \begin{block}{Prueba operativa}Para cada valor usado, registrar quién lo genera, cuándo existe y qué datos intervinieron en su cálculo.\end{block}
\end{frame}

% 21
\begin{frame}{Tres fugas frecuentes}
  \scriptsize
  \begin{tabularx}{\textwidth}{>{\bfseries}p{.2\textwidth}Xp{.31\textwidth}}\toprule
  Tipo & Ejemplo & Control\\\midrule
  Temporal & usar la duración final del bloqueo & cortar variables en $t_0$\\
  Preprocesamiento & normalizar con media de todo el conjunto & ajustar solo con entrenamiento\\
  Duplicados o entidad & misma ventana o vehículo en ambos lados & deduplicar y separar por unidad o grupo\\\bottomrule
  \end{tabularx}
  \begin{block}{Resultado}Quitar una columna sospechosa no basta; la frontera experimental debe aislar toda información futura o compartida.\end{block}
\end{frame}

% 22
\begin{frame}{Generalización}
  \small Generalizar es mantener desempeño en unidades nuevas de la población objetivo.
  \[R(f)=\mathbb{E}_{(X,Y)\sim P_{\mathrm{uso}}}[L(Y,f(X))].\]
  La evaluación estima ese riesgo con datos no usados para ajustar el procedimiento.
  \begin{itemize}
    \item Muestra representativa del uso esperado.
    \item Separación compatible con dependencias temporales y de entidad.
    \item Métrica alineada con el error relevante.
  \end{itemize}
  \begin{block}{Idea}Un valor fuera de muestra es evidencia bajo un escenario de uso definido, no una garantía universal.\end{block}
\end{frame}

% 23
\begin{frame}{Entrenamiento, validación y prueba}
  \small
  \begin{tabularx}{\textwidth}{>{\bfseries}p{.22\textwidth}X}\toprule
  Entrenamiento & ajustar parámetros y transformaciones\\
  Validación & elegir variables, configuraciones y umbral con la métrica principal ya definida\\
  Prueba & medir una vez el procedimiento congelado\\\bottomrule
  \end{tabularx}
  Consultar repetidamente la prueba la convierte en validación.
  \begin{alertblock}{Control}La prueba se reserva para la medición final; cualquier cambio motivado por ella exige una nueva prueba independiente.\end{alertblock}
\end{frame}

% 24
\begin{frame}{Partición aleatoria y estratificada}
  \small
  \begin{itemize}
    \item \concepto{Aleatoria}: asigna unidades al azar; corresponde cuando son intercambiables e independientes.
    \item \concepto{Estratificada}: conserva aproximadamente la proporción de clases en cada partición.
    \item La estratificación estabiliza la medición cuando el bloqueo es poco frecuente.
    \item Ninguna resuelve dependencias entre registros de una entidad ni anticipación temporal.
  \end{itemize}
  \begin{alertblock}{Control}Primero se define qué observaciones pueden separarse; después se conserva la distribución relevante.\end{alertblock}
\end{frame}

% 25
\begin{frame}{Partición por grupos}
  \small Todos los registros de un mismo grupo quedan en una sola partición.
  \begin{center}\begin{tikzpicture}[node distance=8mm,scale=.82,transform shape]
    \node[datum] (t) {Entrenamiento\br A, B, C};
    \node[process,right=of t] (v) {Validación\br D};
    \node[decision,right=of v] (p) {Prueba\br E};
  \end{tikzpicture}\end{center}
  \begin{itemize}
    \item Grupo posible: vehículo, corredor, sensor, cliente o evento.
    \item Evalúa transferencia hacia grupos no observados durante el ajuste.
    \item Evita que patrones propios de una entidad aparezcan a ambos lados.
  \end{itemize}
  \pregunta{¿El uso real predice nuevos instantes de entidades conocidas o entidades completamente nuevas?}
\end{frame}

% 26
\begin{frame}{Partición temporal}
  \small El pasado entrena y el futuro evalúa:
  \begin{center}\begin{tikzpicture}[node distance=2mm,scale=.8,transform shape]
    \node[datum] (e) {Enero--marzo\br entrenamiento}; \node[process,right=of e] (v) {Abril\br validación}; \node[decision,right=of v] (p) {Mayo\br prueba};
    \draw[flow] (e)--(v); \draw[flow] (v)--(p);
  \end{tikzpicture}\end{center}
  \begin{itemize}
    \item Respeta causalidad y simula despliegue prospectivo.
    \item Conserva juntas las ventanas de un mismo incidente cuando corresponde.
    \item Puede revelar cambios de estación, sensores, rutas o comportamiento.
  \end{itemize}
  \begin{alertblock}{Control}Toda variable y transformación se calcula únicamente con información anterior al período evaluado.\end{alertblock}
\end{frame}

% 27
\begin{frame}{Baseline: el mínimo que debe superarse}
  \small Un baseline es una referencia simple, reproducible y disponible con la misma información.
  \begin{itemize}
    \item Clasificación: clase mayoritaria o regla operativa vigente.
    \item Regresión: media o mediana del entrenamiento.
    \item Series temporales: último valor observado o patrón estacional simple.
  \end{itemize}
  \begin{block}{Resultado}Un modelo complejo aporta valor solo si mejora una referencia pertinente en la métrica y el escenario elegidos.\end{block}
\end{frame}

% 28
\begin{frame}{Baseline de clasificación: accuracy engañosa}
  \small En 100 registros hay 20 bloqueos y 80 libres. El baseline siempre predice \texttt{libre}.
  \begin{center}\begin{tabular}{lrr}\toprule
   & Real bloqueado & Real libre\\\midrule
  Predice bloqueado & VP $=0$ & FP $=0$\\ Predice libre & FN $=20$ & VN $=80$\\\bottomrule
  \end{tabular}\end{center}
  \[\mathrm{accuracy}=\frac{0+80}{100}=80\%,\qquad \mathrm{recall}=\frac{0}{0+20}=0\%.\]
  \begin{alertblock}{Lectura}Acierta la clase frecuente, pero no detecta ningún bloqueo; accuracy sola no representa el objetivo.\end{alertblock}
\end{frame}

% 29
\begin{frame}{Baselines de regresión y temporales}
  \scriptsize
  \begin{tabularx}{\textwidth}{p{.28\textwidth}p{.28\textwidth}X}\toprule
  \textbf{Problema} & \textbf{Baseline} & \textbf{Qué controla}\\\midrule
  demanda diaria de agua & mediana del entrenamiento & nivel constante robusto\\
  tiempo de resolución & media del entrenamiento & promedio histórico\\
  velocidad del tramo & último valor observado & persistencia temporal\\
  demanda semanal & valor del mismo día previo & estacionalidad simple\\\bottomrule
  \end{tabularx}
  \begin{block}{Control}Los baselines fijos se estiman con entrenamiento. Un baseline temporal actualiza su estado en cada origen usando solo información disponible hasta ese instante.\end{block}
\end{frame}

% 30
\begin{frame}{Métrica y costo}
  \small La métrica resume errores predictivos; el costo expresa sus consecuencias.
  \begin{tabularx}{\textwidth}{>{\bfseries}p{.18\textwidth}X}\toprule
  FN & predecir libre con bloqueo: ingreso inseguro, detención y demora alta\\
  FP & predecir bloqueado con tramo libre: desvío innecesario\\\bottomrule
  \end{tabularx}
  \begin{itemize}
    \item Un umbral cambia el equilibrio entre FP y FN.
    \item La métrica elegida debe visibilizar el error costoso.
    \item La decisión puede minimizar costo esperado sujeto a restricciones de seguridad.
  \end{itemize}
  \begin{block}{Idea}Ninguna métrica reemplaza la declaración explícita de costos y límites operativos.\end{block}
\end{frame}

% 31
\begin{frame}{Matriz de confusión}
  \small
  \begin{center}\begin{tabular}{lcc}\toprule
   & Real positivo: bloqueado & Real negativo: libre\\\midrule
  Predicción positiva & VP & FP\\ Predicción negativa & FN & VN\\\bottomrule
  \end{tabular}\end{center}
  \begin{itemize}
    \item VP y VN son aciertos; FP es una falsa alarma; FN es un bloqueo omitido.
    \item Total $N=VP+VN+FP+FN$.
    \item Positivos reales $P=VP+FN$; positivos predichos $\hat P=VP+FP$.
  \end{itemize}
  \begin{alertblock}{Control}Definir la clase positiva antes de interpretar cualquier métrica.\end{alertblock}
\end{frame}

% 32
\begin{frame}{Accuracy, precisión, recall y F1}
  \small
  \[\text{accuracy}=\frac{VP+VN}{VP+VN+FP+FN},\qquad \text{precisión}=\frac{VP}{VP+FP},\]
  \[\text{recall}=\frac{VP}{VP+FN},\qquad F_1=2\frac{\text{precisión}\,\text{recall}}{\text{precisión}+\text{recall}}=\frac{2VP}{2VP+FP+FN}.\]
  \begin{itemize}
    \item Precisión: de las alertas emitidas, qué fracción era bloqueo.
    \item Recall: de los bloqueos reales, qué fracción fue detectada.
    \item $F_1$: media armónica de precisión y recall; no incorpora VN.
  \end{itemize}
  \begin{alertblock}{Denominadores}Los cocientes requieren denominador mayor que cero; si no hay positivos reales o predichos, se reporta esa condición.\end{alertblock}
\end{frame}

% 33
\begin{frame}{MAE y RMSE}
  \small Para errores $e_i=y_i-\hat y_i$ en $n$ observaciones:
  \[\mathrm{MAE}=\frac{1}{n}\sum_{i=1}^{n}|e_i|,\qquad \mathrm{RMSE}=\sqrt{\frac{1}{n}\sum_{i=1}^{n}e_i^2}.\]
  \begin{itemize}
    \item Ambos se expresan en la unidad de $y$.
    \item MAE pondera linealmente cada magnitud de error.
    \item RMSE penaliza con mayor fuerza los errores grandes.
  \end{itemize}
  \begin{alertblock}{Control}Comparar ambas métricas sobre las mismas observaciones y declarar si los errores extremos tienen un costo especial.\end{alertblock}
\end{frame}

% 34
\begin{frame}{Validación cruzada K-fold}
  \small Para unidades intercambiables, se divide el conjunto de desarrollo en $K$ folds. Cada fold actúa una vez como validación y los $K-1$ restantes como entrenamiento.
  \[\bar m=\frac{1}{K}\sum_{k=1}^{K}m_k.\]
  \begin{itemize}
    \item Con grupos, cada entidad completa se asigna a un fold.
    \item Con tiempo, se usan ventanas expansivas o deslizantes: solo el pasado entrena y el futuro valida.
    \item La prueba permanece fuera del ciclo completo.
  \end{itemize}
  \begin{block}{Resultado}La estrategia de remuestreo debe reproducir las dependencias y el orden del uso real.\end{block}
\end{frame}

% 35
\begin{frame}{Transformaciones dentro de cada fold}
  \small En cada iteración:
  \begin{enumerate}
    \item ajustar imputación, escalado y codificación con folds de entrenamiento;
    \item transformar entrenamiento y validación con esos valores;
    \item ajustar el modelo en entrenamiento transformado;
    \item medir en validación transformada.
  \end{enumerate}
  \begin{center}\begin{tikzpicture}[node distance=3mm,scale=.69,transform shape]
    \node[datum] (d) {Folds de\br entrenamiento}; \node[process,right=of d] (t) {Ajustar\br transformación};
    \node[process,right=of t] (m) {Ajustar\br modelo}; \node[decision,right=of m] (v) {Aplicar a fold\br de validación};
    \draw[flow] (d)--(t); \draw[flow] (t)--(m); \draw[flow] (m)--(v);
  \end{tikzpicture}\end{center}
  \begin{alertblock}{Control}Ajustar el preprocesamiento con todos los datos filtra información entre folds.\end{alertblock}
\end{frame}

% 36
\begin{frame}{Selección de características}
  \small Elegir variables también aprende de los datos y forma parte del procedimiento.
  \begin{itemize}
    \item \concepto{Disponibilidad}: cada variable debe existir en $t_0$.
    \item \concepto{Pertinencia}: representa una señal compatible con el objetivo.
    \item \concepto{Estabilidad}: significado y medición se sostienen en uso.
    \item \concepto{Selección}: cualquier ranking o descarte basado en datos se ajusta dentro de cada fold.
  \end{itemize}
  \begin{block}{Resultado}Una lista fijada por disponibilidad o conocimiento del dominio puede preespecificarse; usar información de validación o prueba para elegir variables produce optimismo.\end{block}
\end{frame}

% 37
\begin{frame}{Subajuste y sobreajuste}
  \small
  \begin{columns}[T]
    \column{.48\textwidth}
    \centering
    \begin{tikzpicture}[x=1cm,y=1cm,scale=.72]
      \draw[axis] (0,0)--(5.2,0) node[right,font=\scriptsize]{complejidad};
      \draw[axis] (0,0)--(0,3.5) node[above,font=\scriptsize]{error};
      \draw[UAPTeal,very thick] plot[smooth] coordinates{(.3,3)(1.2,2.2)(2.3,1.3)(3.4,.75)(4.8,.4)};
      \draw[UAPRed,very thick] plot[smooth] coordinates{(.3,3.2)(1.2,2.1)(2.4,1.25)(3.2,1.05)(4,1.45)(4.8,2.4)};
      \node[UAPTeal,font=\scriptsize] at (4.1,.45) {entrenamiento};
      \node[UAPRed,font=\scriptsize] at (4.3,2.25) {validación};
      \draw[dashed,UAPGold,thick] (3.2,0)--(3.2,1.05);
    \end{tikzpicture}
    \column{.48\textwidth}
    \scriptsize
    \begin{itemize}
      \item \textbf{Subajuste}: errores altos en entrenamiento y validación; mejorar representación o capacidad.
      \item \textbf{Equilibrio}: validación mínima y brecha controlada; conservar el procedimiento.
      \item \textbf{Sobreajuste}: error de entrenamiento bajo y validación peor; reducir búsqueda, regularizar o reunir datos relevantes.
    \end{itemize}
  \end{columns}
  \begin{block}{Idea}Se elige por desempeño de validación, no por el mejor ajuste al entrenamiento.\end{block}
\end{frame}

% 38
\begin{frame}{Comparación justa y protocolo congelado}
  \small Dos alternativas se comparan con las mismas:
  \begin{itemize}
    \item unidades, folds y período de prueba;
    \item variables disponibles y tratamiento de faltantes;
    \item métrica, clase positiva y regla de agregación;
    \item presupuesto de búsqueda y semillas registradas.
  \end{itemize}
  \begin{block}{Protocolo congelado}Código, versión de datos, particiones, transformaciones, variables, configuración, umbral y métricas quedan fijados antes de abrir la prueba.\end{block}
  \textbf{Resultado:} reproducibilidad es reconstruir el procedimiento y obtener resultados compatibles bajo las mismas condiciones.
\end{frame}

% 39
\begin{frame}{Actividad integradora: diseñar antes de modelar}
  \small Todos los equipos usan una plantilla común y la aplican a su proyecto. El robot o vehículo que estima un bloqueo funciona como caso de referencia compartido.
  \begin{enumerate}
    \item Escriban decisión, salida predictiva y regla de acción por separado.
    \item Definan unidad, población, $t_0$, horizonte, $x$, $y$ y fuente de confirmación.
    \item Marquen tres riesgos de fuga: temporal, preprocesamiento y entidad.
    \item Registren partición, baseline y métrica principal en la plantilla común; justifiquen cada elección por el uso.
  \end{enumerate}
  \begin{block}{Producto de trabajo}Una ficha experimental de una página y un esquema del flujo de datos.\end{block}
\end{frame}

% 40
\begin{frame}{Hito semanal y entregable}
  \scriptsize
  \begin{tabularx}{\textwidth}{>{\bfseries}p{.2\textwidth}X}\toprule
  Sección & Evidencia mínima\\\midrule
  Formulación & decisión, $\hat y$, acción, unidad, población, $t_0$, $H$, $x$, $y$\\
  Partición & diagrama de entrenamiento, validación y prueba con justificación\\
  Controles & fugas identificadas y transformaciones dentro del fold\\
  Evaluación & baseline y métrica principal fijados antes de comparar; costo relevante\\
  Reproducción & versión de datos, semillas y procedimiento congelado\\\bottomrule
  \end{tabularx}
  \begin{block}{Hito}Otro equipo debe poder aplicar los mismos criterios, detectar información indebida y reproducir la evaluación propuesta.\end{block}
\end{frame}

% 41
\begin{frame}{Síntesis}
  \small
  \begin{center}\begin{tikzpicture}[node distance=1.8mm,scale=.62,transform shape]
    \node[decision] (d) {Decisión}; \node[datum,right=of d] (c) {Definición}; \node[datum,right=of c] (p) {Partición};
    \node[process,right=of p] (b) {Baseline}; \node[process,right=of b] (r) {Procedimiento};
    \node[process,right=of r] (v) {Validación}; \node[decision,right=of v] (t) {Prueba final}; \node[decision,right=of t] (a) {Acción};
    \draw[flow] (d)--(c); \draw[flow] (c)--(p); \draw[flow] (p)--(b); \draw[flow] (b)--(r);
    \draw[flow] (r)--(v); \draw[flow] (v)--(t); \draw[flow] (t)--(a);
  \end{tikzpicture}\end{center}
  \begin{itemize}
    \item Aprender estima $\hat y=f(x)$; decidir convierte la salida en acción.
    \item Unidad, población, instante y horizonte definen un ejemplo; la fuga invalida generalización.
    \item La partición representa dependencias y uso; baselines y métricas se conectan con costos.
    \item Preprocesamiento y selección se ajustan dentro de cada fold.
    \item La prueba mide una vez un protocolo reproducible y congelado.
  \end{itemize}
  \begin{block}{Resultado}Un buen experimento hace creíble la conexión entre datos históricos y decisiones futuras.\end{block}
\end{frame}

\end{document}
